\documentclass[11pt,letterpaper]{article}

\usepackage[top=1in, bottom=1in, left=1in, right=1in, headheight=14pt]{geometry}
\usepackage{fontspec}
\setmainfont{DejaVu Serif}
\setsansfont{DejaVu Sans}
\setmonofont{DejaVu Sans Mono}

\usepackage{xcolor}
\usepackage{titlesec}
\usepackage{fancyhdr}
\usepackage{enumitem}
\usepackage{hyperref}
\usepackage{booktabs}
\usepackage{tabularx}
\usepackage{longtable}
\usepackage{colortbl}
\usepackage{multirow}
\usepackage{array}
\usepackage{parskip}
\usepackage{tcolorbox}
\usepackage{graphicx}
\usepackage{lastpage}
\usepackage{amsmath}
\usepackage{amssymb}
\usepackage{textcomp}
\usepackage{siunitx}

\tcbuselibrary{skins,breakable}

% ── Color Scheme: Space / Physics ──────────────────────────────────
\definecolor{primary}{HTML}{311B92}
\definecolor{secondary}{HTML}{0277BD}
\definecolor{tertiary}{HTML}{E65100}
\definecolor{accent}{HTML}{7B1FA2}
\definecolor{lightprimary}{HTML}{EDE7F6}
\definecolor{lightsecondary}{HTML}{E1F5FE}
\definecolor{lighttertiary}{HTML}{FBE9E7}
\definecolor{rowgray}{HTML}{F5F5F5}
\definecolor{bordergray}{HTML}{BDBDBD}
\definecolor{subtlegray}{HTML}{757575}
\definecolor{darktext}{HTML}{212121}

% ── Section Formatting ─────────────────────────────────────────────
\titleformat{\section}
  {\Large\bfseries\sffamily\color{primary}}
  {\thesection}{1em}{}
  [\vspace{-0.5em}\textcolor{bordergray}{\rule{\textwidth}{0.4pt}}]

\titleformat{\subsection}
  {\large\bfseries\sffamily\color{secondary}}
  {\thesubsection}{1em}{}

\titleformat{\subsubsection}
  {\normalsize\bfseries\sffamily\color{tertiary}}
  {\thesubsubsection}{1em}{}

\titlespacing*{\section}{0pt}{1.5em}{0.8em}
\titlespacing*{\subsection}{0pt}{1.2em}{0.5em}
\titlespacing*{\subsubsection}{0pt}{0.8em}{0.3em}

% ── Custom Environments ────────────────────────────────────────────
\newcommand{\stageheader}[2]{%
  \noindent\colorbox{#2}{\makebox[\dimexpr\textwidth-2\fboxsep\relax][l]{%
    \textcolor{white}{\bfseries\sffamily\hspace{6pt}#1\hspace{6pt}}}}%
  \vspace{0.5em}\par%
}

\newtcolorbox{asciibox}{
  colback=rowgray, colframe=bordergray,
  arc=1pt, boxrule=0.5pt,
  left=6pt, right=6pt, top=4pt, bottom=4pt,
  fontupper=\small\ttfamily,
  breakable
}

\newtcolorbox{quotebox}{
  colback=lightprimary, colframe=primary,
  arc=2pt, boxrule=0.5pt,
  left=12pt, right=12pt, top=8pt, bottom=8pt,
  breakable
}

% ── Table Helpers ──────────────────────────────────────────────────
\newcolumntype{L}[1]{>{\raggedright\arraybackslash}p{#1}}
\newcolumntype{C}[1]{>{\centering\arraybackslash}p{#1}}
\newcommand{\tableheader}[1]{\textbf{\sffamily\textcolor{white}{#1}}}
\newcommand{\metafield}[2]{\textbf{\sffamily #1} #2\par}

% ── Headers / Footers ──────────────────────────────────────────────
\pagestyle{fancy}
\fancyhf{}
\fancyhead[L]{\small\sffamily\textcolor{subtlegray}{The Infinite and One Construction Set}}
\fancyhead[R]{\small\sffamily\textcolor{subtlegray}{GSD Mission Package}}
\fancyfoot[C]{\small\sffamily\textcolor{subtlegray}{Page \thepage\ of \pageref{LastPage}}}
\renewcommand{\headrulewidth}{0.4pt}
\renewcommand{\footrulewidth}{0pt}

% ── Hyperlinks ─────────────────────────────────────────────────────
\hypersetup{
  colorlinks=true,
  linkcolor=secondary,
  urlcolor=secondary,
  pdfauthor={GSD Ecosystem --- Tibsfox},
  pdftitle={The Infinite and One Construction Set -- GSD Mission Package},
  pdfsubject={Vision to Mission Pipeline --- 10-Layer Curriculum Architecture},
}

% ── SI unit setup ──────────────────────────────────────────────────
\sisetup{
  per-mode=symbol,
  range-phrase=--,
}

\begin{document}

% ──────────────────────────────────────────────────────────────────
% TITLE PAGE
% ──────────────────────────────────────────────────────────────────
\thispagestyle{empty}
\begin{center}
\vspace*{2.5cm}

{\small\sffamily\textcolor{primary}{\textbf{GSD RESEARCH MISSION PACKAGE --- EXTRA DEEP 10-LAYER}}}

\vspace{0.3cm}
\textcolor{bordergray}{\rule{9cm}{0.4pt}}

\vspace{1.5cm}
{\fontsize{30}{36}\selectfont\bfseries\sffamily\textcolor{primary}{The Infinite and One\\[0.3em]Construction Set}}

\vspace{0.8cm}
{\Large\sffamily\textcolor{secondary}{From Spring Terminal to Hadron Collider:\\A 100-Level Engineering Curriculum}}

\vspace{0.5cm}
{\large\sffamily\itshape\textcolor{tertiary}{Hardware $\rightarrow$ AI $\rightarrow$ Particle Physics $\rightarrow$ Engineering Fidelity}}

\vspace{2.5cm}
\textcolor{bordergray}{\rule{9cm}{0.4pt}}

\vspace{0.8cm}
{\large\sffamily\textcolor{accent}{10 Layers $\cdot$ 100 Levels $\cdot$ 1 Thread}}

\vspace{0.5cm}
{\sffamily\textcolor{accent}{Vision $\rightarrow$ Research $\rightarrow$ Mission}}\\[0.3em]
{\small\sffamily\textcolor{accent}{Full Pipeline Execution}}

\vspace{2.5cm}
\begin{minipage}{0.75\textwidth}
\centering
{\small\sffamily\textcolor{subtlegray}{Date: March 26, 2026 \quad|\quad Status: Ready for GSD Orchestrator}}\\[0.3em]
{\small\sffamily\textcolor{subtlegray}{Activation Profile: Fleet (Full Crew, 20 Roles)}}\\[0.3em]
{\small\sffamily\textcolor{subtlegray}{Estimated: 15--18 context windows across 6--8 sessions}}\\[0.3em]
{\small\sffamily\textcolor{subtlegray}{Codename: CONSTRUCTION-SET-v1.0}}
\end{minipage}
\end{center}
\newpage

% ──────────────────────────────────────────────────────────────────
% TABLE OF CONTENTS
% ──────────────────────────────────────────────────────────────────
\tableofcontents
\newpage

% ==================================================================
%  STAGE 1 — VISION DOCUMENT
% ==================================================================
\stageheader{STAGE 1 --- VISION DOCUMENT}{primary}

\section{The Infinite and One Construction Set --- Vision Guide}

\metafield{Date:}{March 26, 2026}
\metafield{Status:}{Vision --- Ready for Research Phase}
\metafield{Codename:}{CONSTRUCTION-SET-v1.0}
\metafield{Depends On:}{GSD Skill-Creator Core, Amiga Principle Architecture, Silicon Layer Spec}
\metafield{Context:}{10-Layer Fleet Activation, GSD BBS Educational Pack (sibling), GSD Cloud Ops Curriculum (sibling)}

\subsection{Vision}

There is a spring terminal on a piece of cardboard. Push the lever, insert a wire, release.
The spring clamps down with a force of 0.5--0.8 Newton, transmitted through a brass body
alloyed from copper (Z=29) and zinc (Z=30) --- two elements sitting seventeen positions apart
on the periodic table. That terminal costs approximately fifteen cents. Inside the Large Hadron
Collider at CERN, superconducting niobium-titanium magnets --- niobium at position 41, titanium
at position 22 --- bend proton beams traveling at $0.999999991c$ around a 27-kilometer ring.
The spring terminal and the superconducting dipole are the same thing seen at different scales:
\textit{a connection made precise, held by a force calibrated to preserve information through a
mechanical interface}.

The Infinite and One Construction Set is a 100-level curriculum architecture that makes this
continuity visible, walkable, and teachable. ``Infinite'' because the ladder of abstraction has
no ceiling: every element on the periodic table is a universe of solid-state physics; every
logic gate is a window into information theory; every loss curve is a conversation with
probability at the edge of what hardware can compute. ``One'' because a single thread connects
all of it: the faithful transmission of structured information through successive transformations,
with quantifiable degradation at every step. That thread is called \textit{engineering fidelity}.

The curriculum is organized as ten layers of ten levels each. Layer 1 is matter itself: the
periodic table as engineering inventory, the spring terminal as bill of materials made tactile.
Layer 10 is engineering fidelity as meta-principle: the calibration chain that connects a
milliohm measurement on a brass contact to the energy scale of a Higgs boson decay. Every layer
in between is a complete domain in its own right --- electrical signals, digital logic, computing
architecture, software abstraction, network communication, machine intelligence, quantum mechanics,
particle physics --- and every layer is structurally dependent on the one below it.

The GSD ecosystem is the execution environment for this curriculum. The Amiga Principle sits at
its philosophical center. An Amiga 500 with 512 KB of RAM produced graphics, audio, and
multitasking that a contemporary machine with sixteen times the memory could not replicate.
The reason was layered specialization: Agnus handled DMA, Denise handled video, Paula handled
audio, the 68000 handled logic. Each chip did one thing with perfect fidelity; the system's
emergent capability was not the sum of its clock cycles but the product of its architecture.
At Level 100 of this construction set, a student understands exactly why that works --- and can
apply the same principle to an AI inference pipeline, a particle detector readout chain, or a
distributed mesh network humming in a Pacific Northwest barn.

\subsection{Problem Statement}

\begin{enumerate}[leftmargin=*, label=\textbf{P\arabic*.}]

\item \textbf{The Abstraction Ladder is Invisible.} Engineering education teaches layers in
isolation: chemistry, then circuits, then software, then systems. The connections between layers
are implied but never taught. A physics graduate may not know why CMOS logic gates draw zero
static current; a computer science graduate may not know what a transistor physically is. The
resulting knowledge is brittle at every domain boundary --- and domain boundaries are precisely
where interesting engineering happens.

\item \textbf{The Jump from Hobby Electronics to Physics Research is Unmapped.} A curious
sixteen-year-old who has mastered Ohm's Law has no charted path to quantum field theory. Between
a blinking LED and the ATLAS detector at CERN lies decades of curriculum with no coherent
progression. This discontinuity is artificial --- it is the product of institutional silos, not
physical reality. The electron whose charge drives the LED is the same electron whose track is
reconstructed inside the ATLAS silicon strip detector at 40 million events per second.

\item \textbf{Engineering Fidelity is Never Taught as a First Principle.} Every level of the
abstraction hierarchy either preserves or degrades information. A cold solder joint introduces
resistance. A 32-bit floating-point operation loses precision. A noisy training run produces a
miscalibrated model. A poorly terminated cable smears timing resolution in a particle detector.
These are manifestations of a single engineering principle that no curriculum names explicitly:
\textit{information degradation through imperfect transformation}. When you name it, you can
measure it, minimize it, and teach it.

\item \textbf{The Periodic Table is Treated as Chemistry, Not as Hardware.} The 118 elements are
the construction set's complete inventory. Silicon (14) is the semiconductor substrate of
civilization. Copper (29) is the conductor. Tungsten (74) is the gate metal and the filament.
Neodymium (60) is the permanent magnet in every motor and speaker. Niobium (41) is the
superconducting magnet that bends proton beams at CERN. Gold (79) is the corrosion-proof
contact. Students who encounter the periodic table in a chemistry course do not learn it as the
bill of materials for every electronic device ever built. This is a catastrophic pedagogical miss.

\item \textbf{AI Hardware Abstraction Creates Dangerous Fragility.} AI practitioners use GPUs
without understanding what a GPU is, train networks without understanding what a gradient is,
and deploy transformers without understanding what a matrix multiplication costs in watts. This
is not efficiency --- it is fragility. The abstraction will leak. When it does, a practitioner
without foundational knowledge cannot diagnose the failure, cannot estimate its severity, and
cannot propose a principled fix. The path from silicon to intelligence needs to be a continuous
curriculum, not a series of unconnected black boxes.

\end{enumerate}

\subsection{Core Concept}

\textbf{One Thread, Ten Octaves:} Engineering fidelity --- the faithful preservation of information
content through successive transformations --- is the single principle that unifies all ten layers
of the construction set, from a spring terminal's contact resistance to a quark-gluon plasma
reconstruction residual.

\subsubsection{Domain Map: The Ten Layers}

\begin{asciibox}
LEVEL   LAYER           DOMAIN                         FIDELITY UNIT
------  --------------  ----------------------------   ---------------------------
 1- 10  MATTER          Periodic Table / Materials     Resistivity, eV, bandgap
11- 20  SIGNAL          Electrical Foundations         Ohm's Law, SNR, dB
21- 30  LOGIC           Digital / Boolean              Bit error rate, timing slack
31- 40  COMPUTE         Microcontrollers / CPU         Clock cycles, IPC, cache miss
41- 50  SOFTWARE        Compilers / OS / Languages     Cyclomatic complexity, UB rate
51- 60  NETWORK         OSI / TCP-IP / Mesh            Packet loss %, latency ms
61- 70  INTELLIGENCE    ML / Neural Nets / GPU         Loss, accuracy, FLOP/W
71- 80  QUANTUM         QM Foundations / Qubits        Coherence time T2, gate fidelity
81- 90  COLLIDER        Particle Physics / LHC         Luminosity, energy resolution
91-100  FIDELITY        Engineering Meta-Layer         Calibration chain depth, NIST link

Single Binding Thread:
  electron charge -> voltage -> bit -> instruction -> packet ->
  gradient -> qubit -> particle track -> calibration residual ->
  human understanding
\end{asciibox}

\subsubsection{Module Architecture (10-Module Fleet)}

\begin{asciibox}
              THE INFINITE AND ONE CONSTRUCTION SET
                  10-Module Research Architecture

+--M01-MATTER---+  +--M02-SIGNAL---+  +--M03-LOGIC----+
| Periodic Tbl  |  | Ohm / KVL     |  | Boolean /     |
| Elements in   |  | Spring Term.  |  | CMOS / IEEE   |
| Technology    |  | RC / 555      |  | 754 / bfloat  |
+-------+-------+  +-------+-------+  +-------+-------+
        |                  |                  |
        +--------+---------+------------------+
                 |
        +--------v--------+    +--M05-SOFTWARE--+
        |  M04-COMPUTE    +--->| Compilers / OS |
        | MCU / CPU / Mem |    | Rust / UB      |
        +--------+--------+    +------+---------+
                 |                    |
        +--------v--------------------v---------+
        |            M06-NETWORK                |
        | OSI / TCP-IP / Mesh / LHC Timing Net  |
        +--------+------------------------------+
                 |
        +--------v--------------------------+
        |        M07-INTELLIGENCE           |
        | Perceptron -> Transformer / GPU   |
        +--------+---------+---------------+
                 |          |
        +--------v----+  +--v-----------+
        | M08-QUANTUM |  | M09-COLLIDER |
        | QM / Qubits |  | LHC / SM /   |
        | Coherence   |  | AI in HLT    |
        +--------+----+  +----+---------+
                 |             |
        +--------v-------------v--------+
        |         M10-FIDELITY          |
        | Calibration Chain / NIST /    |
        | Meta-Principle Across All 10  |
        +-------------------------------+
\end{asciibox}

\subsection{Research Modules}

\begin{center}
\rowcolors{2}{rowgray}{white}
\begin{tabularx}{\textwidth}{C{0.6cm}L{2.4cm}XC{1.8cm}}
\toprule
\rowcolor{primary}
\tableheader{\#} & \tableheader{Module} & \tableheader{Scope} & \tableheader{Levels} \\
\midrule
M01 & MATTER & Periodic table as engineering inventory; elements in technology; spring terminal anatomy & 1--10 \\
M02 & SIGNAL & Electrical foundations; KVL/KCL; spring terminal experiment; RC circuits; 555 timer & 11--20 \\
M03 & LOGIC & Boolean algebra; 7 gates; CMOS; binary arithmetic; IEEE 754; bfloat16 & 21--30 \\
M04 & COMPUTE & Microcontrollers; CPU architecture; memory hierarchy; Amiga 68000; bus systems & 31--40 \\
M05 & SOFTWARE & Compilers; operating systems; programming languages; Rust ownership; abstraction & 41--50 \\
M06 & NETWORK & OSI model; TCP/IP; mesh networks; LHC timing network; deterministic Ethernet & 51--60 \\
M07 & INTELLIGENCE & ML foundations; perceptron to transformer; GPU silicon; hardware-AI bridge & 61--70 \\
M08 & QUANTUM & Wave-particle duality; Schrodinger equation; qubits; coherence; error correction & 71--80 \\
M09 & COLLIDER & Standard Model; LHC parameters; detector readout; AI in particle physics & 81--90 \\
M10 & FIDELITY & Engineering fidelity as universal meta-principle; calibration chain; traceability & 91--100 \\
\bottomrule
\end{tabularx}
\end{center}

\subsection{Chipset Configuration}

\begin{asciibox}
chipset:
  codename: CONSTRUCTION_SET_v1
  activation_profile: fleet
  parallel_tracks: 3
  wave_count: 6           # W0, W1A, W1B, W1C, W2, W3
  modules:
    - {id: M01, name: matter,        model: sonnet, role: CRAFT-MATTER, tokens: 12000}
    - {id: M02, name: signal,        model: sonnet, role: CRAFT-SIGNAL, tokens: 10000}
    - {id: M03, name: logic,         model: sonnet, role: CRAFT-LOGIC,  tokens: 10000}
    - {id: M04, name: compute,       model: sonnet, role: CRAFT-COMPUTE,tokens: 12000}
    - {id: M05, name: software,      model: sonnet, role: CRAFT-SOFT,   tokens: 10000}
    - {id: M06, name: network,       model: sonnet, role: CRAFT-NET,    tokens:  8000}
    - {id: M07, name: intelligence,  model: opus,   role: CRAFT-AI,     tokens: 16000}
    - {id: M08, name: quantum,       model: opus,   role: CRAFT-QUANT,  tokens: 14000}
    - {id: M09, name: collider,      model: opus,   role: CRAFT-PHYS,   tokens: 16000}
    - {id: M10, name: fidelity,      model: opus,   role: CRAFT-FIDEL,  tokens: 14000}
  orchestrator: {model: opus, role: FLIGHT, cache_ttl: 300s}
  schema_agent:  {model: haiku, role: HAIKU-SCAFFOLD, tokens: 6000}
  model_split:   {opus: 0.30, sonnet: 0.60, haiku: 0.10}
\end{asciibox}

\subsection{Scope Boundaries}

\subsubsection{In Scope (v1.0)}

\begin{itemize}[leftmargin=*]
  \item Ten research modules covering all 100 levels with explicit level-to-level transitions
  \item Periodic table mapping: all 118 elements tabulated, 40+ with engineering applications and device examples
  \item Level 1--20 hands-on experiment sequence: spring terminal cardboard-box progression, complete BOM
  \item Microcontroller-to-CPU architecture progression (L31--L50) including Amiga 68000 case study
  \item Neural network hardware pipeline: GPU silicon to training run, float32 to bfloat16 fidelity
  \item LHC physics primer: full technical parameter set, detector readout chain, AI in High-Level Trigger
  \item Standard Model reference: 17 fundamental particles with mass, charge, role
  \item Engineering fidelity as cross-cutting meta-principle with measurable metrics at all ten layers
  \item GSD skill-creator delivery architecture: each level decomposable into a skill-creator mission component
  \item 12 success criteria, 47-test plan, full verification matrix
\end{itemize}

\subsubsection{Out of Scope (v1.0)}

\begin{itemize}[leftmargin=*]
  \item Full worked curriculum materials (exercises, assessments, grading rubrics) --- v2.0
  \item Quantum computing hardware implementation beyond conceptual primer --- v2.0
  \item Beyond-Standard-Model theories (supersymmetry, string theory, extra dimensions) --- v3.0
  \item Biochemical and biological abstraction layers --- separate mission
  \item Economic and sociological analysis of engineering education equity --- separate mission
  \item Mains-voltage (120V/230V) experiment sequences --- safety boundary, excluded permanently
\end{itemize}

\subsection{Success Criteria}

\begin{enumerate}[leftmargin=*, label=\textbf{SC-\arabic*.}]
  \item The periodic table section maps at least 40 elements to specific engineering applications with documented device examples and primary sources.
  \item The spring terminal experiment sequence provides a complete Level 1--20 hands-on progression: bill of materials, step-by-step procedures, and measurable outcomes per level, all at $\leq$9V.
  \item Every layer transition (L10$\rightarrow$L11, L20$\rightarrow$L21, etc.) is documented with an explicit bridging concept showing how the lower layer enables the upper. Nine bridges total.
  \item The CPU architecture section traces the journey from a single transistor to a full von Neumann machine --- and through the Amiga 68000 as Amiga Principle proof-of-concept --- with no unexplained jumps.
  \item The AI/ML module traces the path from a perceptron to a transformer architecture, with hardware substrate (silicon to data center) described at each milestone.
  \item The LHC section documents the full data acquisition chain from proton-proton collision to reconstructed physics object, including the specific role of graph neural networks in the High-Level Trigger.
  \item Engineering fidelity metrics are defined for each of the ten layers, producing a single calibration chain from spring terminal contact resistance to LHC detector energy resolution, with NIST traceability noted.
  \item The quantum mechanics module bridges classical probability to quantum superposition without assuming prior quantum knowledge, achieving conceptual continuity from Layer 3 (logic) to Layer 8.
  \item All quantitative claims (physical constants, detector specs, transistor counts, training compute, GPU specs) are attributed to specific primary technical sources.
  \item The research reference is self-contained for a reader with a high-school physics background: no undefined terms, no unexplained references, complete bibliography.
  \item At least three explicit inter-layer dependency relationships are traced with concrete physical examples connecting non-adjacent layers (e.g., periodic table $\rightarrow$ LHC superconductor).
  \item The curriculum architecture supports GSD skill-creator delivery: each level is decomposable into a skill-creator mission component with defined input, output, and fidelity metric.
\end{enumerate}

\subsection{Relationship to Other Documents}

\begin{center}
\rowcolors{2}{rowgray}{white}
\begin{tabularx}{\textwidth}{L{3.8cm}L{2cm}X}
\toprule
\rowcolor{primary}
\tableheader{Document} & \tableheader{Rel.} & \tableheader{Notes} \\
\midrule
GSD Skill-Creator Core & Execution & This curriculum runs as skill-creator missions; each level is a potential skill \\
GSD Amiga Vision & Philosophy & Amiga Principle is the curriculum's central proof-of-concept \\
The Space Between (TSB) & Foundational & Mathematical spine; fidelity as information-theoretic concept \\
GSD Silicon Layer Spec & Sibling & Hardware substrate; companion document to Levels 31--50 \\
GSD BBS Educational Pack & Sibling & Delivery layer; BBS-style level progression for L1--30 \\
GSD Cloud Ops Curriculum & Sibling & Overlap at network/distributed layer (Levels 51--60) \\
GSD Instruction Set Arch. & Sibling & CPU/ISA depth for Levels 41--50 \\
Fox Infrastructure Group & Infra & Physical cluster for curriculum execution; mesh spec relevant to L51--60 \\
GSD Mesh Prototype Spec & Infra & 10-node cluster as curriculum lab substrate \\
\bottomrule
\end{tabularx}
\end{center}

\subsection{The Through-Line}

The title carries a double meaning, and both are deliberate. ``Infinite'' refers to the depth
available at any point in the construction set. Pick any element on the periodic table and
follow it: silicon at position 14 leads to semiconductor physics, to quantum mechanics, to the
band theory of solids, to CMOS fabrication, to Moore's Law, to the GPU running the model that
is reading this text. Every node in the construction set is a universe if you go deep enough.
``One'' refers to the single thread that connects all of them: information, and whether it
survives the passage from one level to the next.

In the GSD ecosystem this is the space between --- the meaning that lives in the connections, not
in the nodes. The spring terminal's spring is not interesting because it is a spring; it is
interesting because it exerts a force that maintains electrical continuity between two conductors,
preserving signal fidelity across a mechanical interface. The Large Hadron Collider is not
interesting because it is large; it is interesting because it preserves the identity of a proton's
constituent quarks across a 26.7-kilometer path bent by superconducting fields, so that when two
protons annihilate at $\sqrt{s} = 13.6$ TeV, the resulting particle shower can be reconstructed to
millimeter spatial precision and traced back to a quantum event that lasted $10^{-25}$ seconds.

The spring and the superconducting dipole are doing the same job at different scales. This
curriculum makes that continuity visible. It is the Amiga Principle applied to all of physics:
architectural elegance over brute force, from the electron to the quark and back.

\newpage

% ==================================================================
%  STAGE 2 — RESEARCH REFERENCE
% ==================================================================
\stageheader{STAGE 2 --- RESEARCH REFERENCE}{secondary}

\section{The Infinite and One Construction Set --- Research Reference}

\subsection{How to Use This Document}

This reference provides the technical substrate for all ten curriculum modules. Each module
section contains foundational concepts, key quantitative data, engineering parameters, and
cross-references to adjacent modules. It is designed for skill-creator agents producing curriculum
content at each of the 100 levels, and for human instructors preparing sessions.

\textbf{Key source organizations for this research domain:}
\begin{itemize}[leftmargin=*]
  \item \textbf{CERN} --- LHC technical documentation, Run 3 accelerator parameters, detector papers
  \item \textbf{NIST} --- Atomic weights, resistivity values, CODATA fundamental constants
  \item \textbf{Particle Data Group (PDG/LBNL)} --- Standard Model particle masses and properties
  \item \textbf{IEEE} --- Circuit and signal standards; GPU architecture; wireless protocols
  \item \textbf{ACM} --- Computing curricula; software engineering body of knowledge
  \item \textbf{Royal Society of Chemistry} --- Periodic table element data and applications
  \item \textbf{MIT OpenCourseWare} --- Curriculum architecture models (6.004, 6.012, 8.05)
  \item \textbf{arXiv.org} --- Preprint physics and ML research; GNN-for-tracking papers
  \item \textbf{NVIDIA} --- GPU architecture white papers; CUDA programming model
\end{itemize}

% ─────────────────────────────────────────────────────────────────
\subsection{M01 --- MATTER: The Periodic Table as Engineering Inventory}
% ─────────────────────────────────────────────────────────────────

\subsubsection{The Construction Set's Bill of Materials}

The periodic table, arranged by atomic number $Z$ from hydrogen ($Z=1$) to oganesson ($Z=118$),
encodes the complete inventory of matter available for physical construction. Mendeleev published
the first periodic table in 1869; the modern quantum-mechanical explanation (electron shells and
orbital substructure) was completed in the 1920s. For an engineer, the relevant atomic properties
are: electrical resistivity, thermal conductivity, melting point, magnetic permeability, atomic
radius, and electronegativity (which governs alloy behavior and chemical bonding at interfaces).

\textbf{Conductors --- Group 11 and neighbors:}
Copper (Cu, $Z=29$) has resistivity $\rho = 1.68 \times 10^{-8}\ \Omega\cdot\text{m}$ at
$20^\circ$C, making it the primary conductor in virtually all electronic wiring. Silver (Ag, $Z=47$)
has $\rho = 1.59 \times 10^{-8}\ \Omega\cdot\text{m}$ but is economically prohibitive at scale. Gold
(Au, $Z=79$) is used for connector contacts because its oxide is conductive --- unlike copper and
silver, gold does not form an insulating oxide layer, maintaining fidelity at the interface.
Aluminum (Al, $Z=13$) is used in high-voltage transmission lines for its strength-to-weight ratio
despite $\rho = 2.82 \times 10^{-8}\ \Omega\cdot\text{m}$.

\textbf{Semiconductors --- Group 14:}
Silicon (Si, $Z=14$) is the substrate of modern computing, with bandgap $E_g = 1.12$ eV at room
temperature. Doped with phosphorus (P, $Z=15$), it becomes n-type (extra electrons); doped with
boron (B, $Z=5$), p-type (extra holes). Germanium (Ge, $Z=32$), bandgap $0.66$ eV, was the
original transistor material (Shockley, Bardeen, Brattain, 1947) before silicon displaced it.
Gallium arsenide (GaAs, compound of Ga $Z=31$ and As $Z=33$) has higher electron mobility than
Si, enabling high-frequency RF devices and III--V solar cells. Silicon carbide (SiC, Si + C $Z=6$),
bandgap $3.26$ eV, enables power electronics above 1 kV.

\textbf{Superconductors --- the LHC connection:}
Below a critical temperature $T_c$, certain materials exhibit zero electrical resistance.
The LHC's 1,232 main bending dipole magnets use niobium-titanium alloy (Nb$_{0.47}$Ti$_{0.53}$),
operating at 1.9 K in superfluid helium. Niobium ($Z=41$, $T_c = 9.25$ K) combined with titanium
($Z=22$) forms a ductile alloy with $T_c \approx 9$ K and mechanical properties suitable for magnet
coil winding. Each dipole carries 11,850 A producing a field of 8.33 T (Run 2) / 8.60 T (Run 3).
This is the periodic table completing its arc: from a spring terminal's brass body (Cu + Zn) to
a superconducting coil (Nb + Ti), the same construction set, at scales separated by 17 orders of
magnitude in energy.

\begin{center}
\rowcolors{2}{rowgray}{white}
\begin{tabularx}{\textwidth}{L{1.8cm}C{0.7cm}XL{3cm}}
\toprule
\rowcolor{primary}
\tableheader{Element} & \tableheader{Z} & \tableheader{Primary Engineering Role} & \tableheader{Device Example} \\
\midrule
Hydrogen & 1 & Fuel cells, coolant & PEM fuel cells \\
Carbon & 6 & Resistors, nanotubes, graphene & Carbon film resistors \\
Boron & 5 & P-type silicon dopant & CMOS PMOS transistors \\
Silicon & 14 & Semiconductor substrate & CPUs, DRAM, photovoltaics \\
Phosphorus & 15 & N-type silicon dopant & CMOS NMOS transistors \\
Aluminum & 13 & Conductor, IC interconnect & PCB traces, DRAM interconnect \\
Copper & 29 & Primary conductor ($\rho$ = 1.68e-8) & Wire, PCB layers, motor windings \\
Zinc & 30 & Brass alloy component & WAGO/spring terminal bodies \\
Nickel & 28 & Alloy, corrosion barrier & Connector nickel underplate \\
Germanium & 32 & Legacy semiconductor & First transistors (1947) \\
Gallium & 31 & GaAs, LEDs, InGaN & RF ICs, LED phosphors \\
Arsenic & 33 & GaAs, n-type III-V dopant & High-speed RF ICs \\
Silver & 47 & Low-$\rho$ conductor, RFID & RFID antennas, solar busbars \\
Indium & 49 & ITO transparent conductor & LCD/OLED displays \\
Tin & 50 & Solder (SAC305: Sn-Ag-Cu) & PCB assembly joints \\
Neodymium & 60 & NdFeB permanent magnets & Hard drives, motors, speakers \\
Tantalum & 73 & Capacitor dielectric & Tantalum electrolytic caps \\
Tungsten & 74 & Gate metal, filament, X-ray anode & IC gates, light bulbs, X-ray tubes \\
Gold & 79 & Corrosion-proof contacts & IC bond wires, edge connector pins \\
Niobium & 41 & Superconducting magnets & LHC dipole coils \\
Titanium & 22 & NbTi alloy, structural aerospace & LHC coils, RF cavity support \\
\bottomrule
\end{tabularx}
\end{center}

\subsubsection{Spring Terminal Anatomy: A Periodic Table Case Study (Levels 1--5)}

A lever-action spring terminal (WAGO 221 series or equivalent) contains: brass body
(70\% Cu + 30\% Zn, chosen for machinability and spring-back coefficient), stainless steel
spring (Fe + Cr + Ni: strength, corrosion resistance), and tinned copper bus bar. Clamping
force is calibrated to 0.5--0.8 N for AWG 28 wire; contact resistance $< 0.5$ m$\Omega$.
Voltage rating: 600 V. Current rating: 20 A. This fifteen-cent component is a precision
fidelity instrument: it maintains $< 0.5$ m$\Omega$ contact resistance across 10,000
insertion-removal cycles. At Level 1, a student names the elements in the terminal. At Level 5,
the student measures contact resistance with a milliohm meter and traces the measurement to a
NIST-calibrated resistance standard. The spring terminal is the curriculum's first calibration
chain.

% ─────────────────────────────────────────────────────────────────
\subsection{M02 --- SIGNAL: Electrical Foundations}
% ─────────────────────────────────────────────────────────────────

\subsubsection{The Cardboard Box Experiment (Levels 11--15)}

The canonical Level 11 experiment: a 9 V battery, a 470 $\Omega$ carbon film resistor, a red LED,
and two spring terminals on a cardboard base. Ohm's Law governs: $V = IR$. With 9 V supply and
LED forward voltage $V_f = 1.8$ V, the resistor sees 7.2 V. Current: $I = 7.2 / 470 \approx
15.3$ mA. Power dissipated in resistor: $P = I^2 R = 0.110$ W (well within 0.25 W rating).
Power delivered to LED: $P = IV_f = 0.028$ W. This is engineering fidelity at its most
elemental: the intent (``illuminate the LED'') is faithfully transmitted through the material
stack (battery $\rightarrow$ conductor $\rightarrow$ resistor $\rightarrow$ LED) with losses
quantified at every node.

\textbf{Kirchhoff's Laws (Level 12--13):}
\begin{itemize}[leftmargin=*]
  \item \textbf{KCL} (Current Law): $\sum I_{\text{in}} = \sum I_{\text{out}}$ at any node.
    Conservation of charge --- a direct expression of $\nabla \cdot \mathbf{J} = -\partial\rho/\partial t$.
  \item \textbf{KVL} (Voltage Law): $\sum V_k = 0$ around any closed loop. Conservation of energy.
\end{itemize}
These two laws are sufficient to solve any linear DC circuit. They are the signal layer's
encoding of physical conservation laws --- the fidelity constraint at Layer 2.

\subsubsection{Reactive Components and Frequency (Levels 16--19)}

Capacitors ($C$, farads) store charge: $Q = CV$; impedance $Z_C = 1/(j\omega C)$. Inductors
($L$, henrys) store magnetic flux: $V = L \cdot dI/dt$; impedance $Z_L = j\omega L$. Together
they resonate at $f_0 = 1/(2\pi\sqrt{LC})$. The 555 timer IC (Hans Camenzind, 1971; Texas
Instruments, 1972; still in production) uses an RC network for timing: in astable mode,
$t_\text{high} = 0.693(R_A + R_B)C$, $t_\text{low} = 0.693 R_B C$. At Level 18, a student
changes $R_B$ and observes a changed blink rate: the first parameterized, measurable
software-physical loop in the curriculum.

\subsubsection{Signal Integrity as Fidelity Metric (Level 20)}

At high frequencies, transmission lines exhibit characteristic impedance
$Z_0 = \sqrt{L^\prime / C^\prime}$ per unit length ($L^\prime$, $C^\prime$ are distributed
inductance and capacitance). A 50 $\Omega$ coaxial cable must be terminated with a 50 $\Omega$
load; mismatch produces reflections with reflection coefficient
$\Gamma = (Z_L - Z_0)/(Z_L + Z_0)$. Open-circuit termination gives $\Gamma = +1$ (full
positive reflection); short-circuit $\Gamma = -1$ (full negative reflection). Level 20 is the
curriculum's first encounter with impedance mismatch as a fidelity failure --- the same principle
that operates at the LHC, where 50 $\Omega$ coaxial signal cables connect detector front-end
electronics to trigger boards over runs up to 100 m.

% ─────────────────────────────────────────────────────────────────
\subsection{M03 --- LOGIC: Digital Foundations}
% ─────────────────────────────────────────────────────────────────

\subsubsection{Boolean Algebra and the Seven Gates (Levels 21--25)}

George Boole (1854, \textit{An Investigation of the Laws of Thought}) formalized binary logic;
Claude Shannon (1937, MIT master's thesis) proved equivalence to switching circuits. The seven
fundamental gates --- AND, OR, NOT, NAND, NOR, XOR, XNOR --- and their truth tables form the
complete vocabulary of combinational logic. NAND is functionally complete: any Boolean function
can be expressed using only NAND gates. This is the curriculum's first encounter with minimalism
as architectural principle: the Amiga Principle at the logic layer.

\textbf{CMOS Technology (Level 25--27):}
Complementary Metal-Oxide-Semiconductor pairs NMOS (n-channel) and PMOS (p-channel) transistors.
A CMOS inverter draws zero static current: when input is HIGH, NMOS conducts, PMOS cuts off; when
LOW, PMOS conducts, NMOS cuts off. Current flows only during switching transitions
($P \propto \alpha C_{load} V_{DD}^2 f$). A 7 nm CMOS process node (2021--2025) achieves
$\sim$100 million transistors per mm$^2$. The Apple M4 (2024) contains 28 billion transistors.
The fidelity lesson: CMOS power consumption scales with frequency, not static state. This is why
power management is a first-order concern in every layer from the microcontroller to the LHC.

\subsubsection{Binary Arithmetic and Floating-Point Fidelity (Levels 28--30)}

Two's complement signed integers: $n$-bit range is $[-2^{n-1}, 2^{n-1}-1]$. IEEE 754 float32:
1 sign bit, 8 exponent bits, 23 mantissa bits; range $\pm 3.4 \times 10^{38}$; machine epsilon
$\varepsilon \approx 1.19 \times 10^{-7}$ (approximately 7 significant decimal digits). IEEE 754
float64 (double precision): 1 sign, 11 exponent, 52 mantissa; $\varepsilon \approx 2.22 \times
10^{-16}$. Google bfloat16 (brain floating point, 2018): 1 sign, 8 exponent, 7 mantissa; same
exponent range as float32 but $4\times$ lower memory cost. The fidelity implication: every
floating-point operation introduces rounding error. A neural network with 100 layers of float32
matrix multiplications can accumulate gradient errors that exceed the signal. bfloat16 was
designed specifically for AI training: preserve exponent range (critical for gradient magnitude
diversity) while halving memory bandwidth. This is precision engineering at the numeric format
layer.

% ─────────────────────────────────────────────────────────────────
\subsection{M04 --- COMPUTE: Microcontrollers to CPU Architecture}
% ─────────────────────────────────────────────────────────────────

\subsubsection{The Microcontroller as First Complete System (Levels 31--40)}

An Arduino Uno contains an ATmega328P: 8-bit AVR RISC processor, 16 MHz, 5 V; 32 KB flash
(program storage), 2 KB SRAM, 1 KB EEPROM; 14 digital I/O pins; 6 analog inputs (10-bit ADC,
0--5 V, 1,024 steps; LSB = 4.9 mV); 6 PWM outputs (8-bit, 490 Hz default); UART, SPI, I2C.
At Level 31, a student blinks an LED from software (``Hello, World'' of hardware). At Level 40,
the student reads an analog sensor, applies a calibration polynomial
($T = a_0 + a_1 V + a_2 V^2$, coefficients from datasheet), transmits calibrated temperature
over I2C to an OLED display. This is the curriculum's first complete fidelity loop: a physical
measurand (temperature) converts to an electrical signal, digitizes to a 10-bit integer, processes
through an algorithm, and expresses as a human-readable number with quantified uncertainty.

\subsubsection{CPU Architecture: The Amiga Case Study (Levels 41--50)}

The five-stage RISC pipeline (Hennessy and Patterson, 1990): IF (Instruction Fetch), ID
(Instruction Decode), EX (Execute), MEM (Memory Access), WB (Write Back). With perfect pipelining,
one instruction completes per clock cycle (CPI = 1.0). Data hazards require stalls or forwarding;
control hazards (branches) require prediction. A modern out-of-order superscalar processor
(AMD Zen 4, 2022) sustains $> 5$ instructions per clock (IPC) through dynamic scheduling,
register renaming, and speculative execution.

\textbf{The Amiga 68000 and the Amiga Principle:}
The Motorola MC68000 (1979) is a 16/32-bit CISC processor (external 16-bit data bus, internal
32-bit registers, 24-bit address bus; 16 MHz maximum in the Amiga 500). But it is not the CPU
that makes the Amiga remarkable: it is the chipset. Agnus ($\rightarrow$ Fat Agnus $\rightarrow$
Alice): DMA controller, chip memory arbitrator, horizontal/vertical beam position registers.
Denise ($\rightarrow$ Lisa): bitplane display, sprite hardware, color index lookup. Paula: 4-channel
PCM audio DMA, 1-bit serial I/O, disk DMA. The 68000 is freed from all I/O overhead, handling
only application logic. Result: a 7.16 MHz Amiga 500 (1987, \$595) produced 4096-color HAM
graphics and hardware-mixed 4-channel audio simultaneously, while an IBM AT at 16 MHz with EGA
(64 colors) could not approach it. The lesson for every layer above: clock speed is not fidelity.
Architectural leverage is fidelity.

\subsubsection{Memory Hierarchy (Level 48--50)}

\begin{center}
\rowcolors{2}{rowgray}{white}
\begin{tabularx}{\textwidth}{L{2.3cm}L{2cm}L{2.8cm}X}
\toprule
\rowcolor{primary}
\tableheader{Level} & \tableheader{Technology} & \tableheader{Latency (2025)} & \tableheader{Capacity (typical)} \\
\midrule
L1 Cache & SRAM & 1--4 cycles & 32--64 KB per core \\
L2 Cache & SRAM & 4--14 cycles & 256 KB -- 1 MB \\
L3 Cache & SRAM & 20--60 cycles & 8--128 MB shared \\
Main RAM & DRAM & 60--100 ns & 16--256 GB DDR5 \\
NVMe SSD & NAND Flash & 50--100 $\mu$s & 1--8 TB \\
HDD & Magnetic & 5--15 ms & 4--20 TB \\
\bottomrule
\end{tabularx}
\end{center}

Each level in the memory hierarchy is a fidelity trade-off: lower latency costs more money and
power per bit. Cache miss rate is the fidelity metric: the fraction of accesses that must proceed
to a slower tier.

% ─────────────────────────────────────────────────────────────────
\subsection{M05 --- SOFTWARE: Abstraction Layers}
% ─────────────────────────────────────────────────────────────────

\subsubsection{Compiler Pipeline and Undefined Behavior (Levels 41--47)}

A C compiler (clang/LLVM) performs: lexical analysis $\rightarrow$ parsing (AST) $\rightarrow$ 
semantic analysis $\rightarrow$ IR generation (LLVM IR) $\rightarrow$ optimization passes
(inlining, vectorization, constant folding, dead code elimination) $\rightarrow$ code generation
$\rightarrow$ linking. Optimization passes exploit the C standard's Undefined Behavior (UB) rules:
signed integer overflow is UB; the compiler may assume it does not occur and optimize accordingly,
producing programs that are correct in testing and incorrect in production. This is the software
layer's fidelity failure mode: the contract between programmer intent and machine execution is
silently violated when UB is present. Measuring UB density (e.g., via UBSan instrumentation) is
the software layer's fidelity metric.

\subsubsection{Operating Systems as Isolation Machines (Levels 48--52)}

A kernel's primary function is to multiplex hardware resources across processes while enforcing
memory isolation. Virtual memory maps every process's address space to physical pages; the MMU
(Memory Management Unit) in the CPU performs this translation on every load/store, with a
Translation Lookaside Buffer (TLB) caching recent translations ($\sim$1--4 cycle hit,
$\sim$50--100 cycle miss). Linux runs on a Raspberry Pi 4 (Cortex-A72, 1.8 GHz, 4 GB LPDDR4)
and on an AMD EPYC 9654 (96 cores, 2 TB RAM) with the same POSIX interface. The fidelity of the
OS abstraction is measured by the ratio of application cycles to total cycles (eliminating
overhead from context switches, system calls, and kernel interrupts).

\subsubsection{Rust and Structural Memory Safety (Levels 53--55)}

Rust's ownership and borrow-checker system enforces memory safety at compile time: no
use-after-free, no double-free, no null pointer dereference, no data races --- all provable
without garbage collection, without runtime overhead. The GSD skill-creator uses Rust for its
performance-critical components. This is the software layer's instance of the curriculum's deepest
principle: \textit{safety by structure, not by policy}. A program that cannot compile an unsafe
memory access is architecturally precluded from that failure mode. Engineering fidelity made
structural.

% ─────────────────────────────────────────────────────────────────
\subsection{M06 --- NETWORK: Communications Architecture}
% ─────────────────────────────────────────────────────────────────

\subsubsection{The OSI Model as a Fidelity Stack (Levels 51--58)}

The seven-layer OSI model is the clearest expression of the curriculum's layered fidelity concept
outside of physics itself:

\begin{center}
\rowcolors{2}{rowgray}{white}
\begin{tabularx}{\textwidth}{C{0.6cm}L{2.4cm}XL{2.8cm}}
\toprule
\rowcolor{primary}
\tableheader{L} & \tableheader{Layer} & \tableheader{Fidelity Function} & \tableheader{Protocol Examples} \\
\midrule
7 & Application & End-to-end semantic fidelity & HTTP/3, DNS, MQTT, gRPC \\
6 & Presentation & Data format translation & TLS 1.3, Protocol Buffers \\
5 & Session & Connection state management & WebSockets, RPC sessions \\
4 & Transport & Reliable delivery guarantee & TCP, UDP, QUIC \\
3 & Network & Addressing and routing & IPv4/6, ICMP, BGP \\
2 & Data Link & Frame-level error detection & Ethernet, 802.11 WiFi \\
1 & Physical & Signal-level fidelity & Copper, fibre, RF, light \\
\bottomrule
\end{tabularx}
\end{center}

Each layer receives imperfect fidelity from below and presents improved fidelity above. The
Physical layer deals with attenuation and thermal noise; the Data Link layer adds CRC-32 error
detection (Ethernet: Hamming distance 4, detects all 1-, 2-, 3-bit errors in a frame); the
Transport layer adds retransmission, congestion control, and flow control (TCP). This is the
network layer's encoding of the curriculum's central principle: add fidelity at each level by
engineering the failure modes of the level below.

\subsubsection{Mesh Networks and the LHC Timing Network (Levels 58--60)}

The GSD mesh cluster (10-node, Agnus/Denise/Paula/68000 node types) uses a mesh topology for
link redundancy. The LHC's accelerator control and timing network provides the extreme-end
comparison: 600,000 accelerator control channels are synchronized using IEEE 802.1Qbv
Time-Sensitive Networking (TSN) to sub-microsecond timing. The Timing, Trigger, and Control (TTC)
system distributes a 40.079 MHz clock derived from the LHC RF system to all detector front-end
electronics, achieving $< 1$ ns timing resolution across 27 km. This is network fidelity at its
most demanding: a single missed clock edge corrupts an event timestamp, breaking the causal chain
from collision to physics result.

% ─────────────────────────────────────────────────────────────────
\subsection{M07 --- INTELLIGENCE: AI and Machine Learning}
% ─────────────────────────────────────────────────────────────────

\subsubsection{The Hardware-AI Bridge: Perceptron to Transformer (Levels 61--70)}

Frank Rosenblatt's perceptron (1958) computes $\hat{y} = \text{sgn}(\mathbf{w} \cdot \mathbf{x} + b)$
with update rule $\mathbf{w} \leftarrow \mathbf{w} + \eta(y - \hat{y})\mathbf{x}$. This is
computable on any CPU. A 2024-era transformer (Llama 3.1 405B: 405 billion parameters) requires
approximately $6 \times 10^{24}$ FLOPs to train once. The hardware substrate determines what is
\textit{possible} at each era:

\begin{center}
\rowcolors{2}{rowgray}{white}
\begin{tabularx}{\textwidth}{L{1.5cm}L{3.5cm}XL{2.5cm}}
\toprule
\rowcolor{primary}
\tableheader{Level} & \tableheader{Architecture} & \tableheader{Hardware Substrate} & \tableheader{Fidelity Metric} \\
\midrule
L61 & Perceptron & Any CPU / Arduino Uno & Binary accuracy \\
L63 & MLP (3-layer) & Desktop CPU, numpy & MSE / cross-entropy \\
L65 & CNN (ResNet-18) & Raspberry Pi + USB GPU & Top-1 accuracy \\
L67 & RNN / LSTM & Single GPU (RTX 3080) & Perplexity \\
L69 & Transformer (GPT-2) & 1--4 GPU workstation & BLEU / perplexity \\
L70 & Large LLM (70B+) & GPU cluster / TPU pod & Benchmark suites \\
\bottomrule
\end{tabularx}
\end{center}

\subsubsection{GPU Architecture: The AI Spring Terminal (Level 67--70)}

An NVIDIA RTX 5090 (2025) contains 21,760 CUDA cores, 680 4th-generation Tensor Cores,
32 GB GDDR7 at 1.79 TB/s bandwidth. FP32 peak: 109.8 TFLOPS; bfloat16 Tensor Core peak:
$\approx$1,740 TFLOPS. A matrix multiplication of two $4096 \times 4096$ matrices requires
$2 \times 4096^3 = 1.37 \times 10^{11}$ FLOPs: completed in $\approx$0.8 ms on an RTX 5090
Tensor Core, versus $\sim$137 s on a single-threaded CPU core at 1 GFLOPS. The GPU is the
spring terminal of AI: a specialized connector between the mathematical abstraction of a neural
network (floating-point matrix operations) and the physical silicon that executes it.

\subsubsection{Neural Network Fidelity: Loss as Engineering Signal (Level 64--66)}

The loss function is the fidelity metric of a neural network. Cross-entropy loss:
$\mathcal{L} = -\sum_i y_i \log \hat{y}_i$. Gradient descent: $\mathbf{w} \leftarrow \mathbf{w}
- \eta \nabla_\mathbf{w}\mathcal{L}$. Gradient vanishing in deep networks (gradients shrink
exponentially with depth via chain rule) is an engineering fidelity failure: the training signal
does not propagate faithfully through 100+ layers. Solutions are fidelity engineering: residual
connections (He et al., ResNet 2016) provide identity bypass paths; layer normalization
(transformers) re-centers gradient distributions; He/Glorot initialization sets initial weight
scales so that activation variance is preserved forward and gradient variance is preserved
backward. These are not tricks --- they are fidelity engineering at the machine-learning layer,
applying the same principle as impedance matching at the signal layer.

% ─────────────────────────────────────────────────────────────────
\subsection{M08 --- QUANTUM: Foundations of Quantum Mechanics}
% ─────────────────────────────────────────────────────────────────

\subsubsection{Wave-Particle Duality as the Ultimate Fidelity Problem (Levels 71--75)}

The double-slit experiment (Jönsson, 1961 for electrons; Arndt et al., 1999 for $C_{60}$
fullerenes) demonstrates that a single particle produces an interference pattern --- implying
it traverses both paths simultaneously. When a detector is placed at either slit, the pattern
disappears: measurement collapses the quantum state. This is the deepest statement of the
fidelity problem in all of physics: at the quantum layer, \textit{observation changes the
observed}. Classical fidelity --- transmit without altering --- is physically impossible at
this scale. Information and measurement are entangled; you cannot extract information without
disturbing the system.

\textbf{The Schr\"{o}dinger Equation:}
\[i\hbar \frac{\partial \psi}{\partial t} = \hat{H}\psi\]
where $\psi(\mathbf{r},t)$ is the complex-valued wave function, $\hbar = 1.055 \times 10^{-34}$
J$\cdot$s is the reduced Planck constant (NIST CODATA 2018), and $\hat{H}$ is the Hamiltonian
operator (total energy). The probability density of finding the particle at $\mathbf{r}$ is
$|\psi(\mathbf{r},t)|^2$.

\subsubsection{Qubits and Quantum Coherence (Levels 76--80)}

A classical bit is in state $\{0,1\}$. A qubit is in superposition: $|\psi\rangle = \alpha|0\rangle
+ \beta|1\rangle$ where $|\alpha|^2 + |\beta|^2 = 1$ (Born rule). Two qubits can be entangled:
$|\Phi^+\rangle = (|00\rangle + |11\rangle)/\sqrt{2}$; measuring one qubit instantly determines
the other's state regardless of spatial separation (Bell, 1964; Aspect et al., 1982, 2022 Nobel Prize).

\textbf{Coherence time $T_2$} is the fidelity metric of a quantum processor: the duration over
which a qubit maintains phase coherence before coupling to the environment (decoherence) destroys
the superposition. Recent hardware:

\begin{center}
\rowcolors{2}{rowgray}{white}
\begin{tabularx}{\textwidth}{L{3cm}C{1.5cm}L{2.5cm}XC{1.8cm}}
\toprule
\rowcolor{primary}
\tableheader{System} & \tableheader{Qubits} & \tableheader{Year} & \tableheader{Technology} & \tableheader{$T_2$ ($\mu$s)} \\
\midrule
IBM Falcon r4 & 27 & 2021 & Superconducting & $\sim$100 \\
IBM Eagle & 127 & 2021 & Superconducting & $\sim$100 \\
IBM Heron r2 & 133 & 2024 & Superconducting & $\sim$150 \\
Google Willow & 105 & 2024 & Superconducting & $\sim$100 \\
IonQ Forte & 36 & 2024 & Trapped ion & $> 10^6$ \\
Quantinuum H2 & 56 & 2024 & Trapped ion & $> 10^6$ \\
\bottomrule
\end{tabularx}
\end{center}

Google Willow (2024) demonstrated below-threshold quantum error correction: adding qubits
\textit{reduces} logical error rate, a prerequisite for fault-tolerant quantum computing. The
analogy to Layer 2: just as ECC DRAM corrects single-bit errors by adding redundant bits, quantum
error correction corrects qubit errors by encoding one logical qubit in many physical qubits.
Fidelity through redundancy, at the quantum layer.

% ─────────────────────────────────────────────────────────────────
\subsection{M09 --- COLLIDER: Particle Physics and the LHC}
% ─────────────────────────────────────────────────────────────────

\subsubsection{The Standard Model: The Periodic Table's Successor (Levels 81--85)}

The Standard Model of particle physics, completed $\sim$1973 and confirmed through 2012 (Higgs
boson discovery), classifies all known fundamental particles and their interactions. It is the
periodic table's successor: the elements are not the fundamental building blocks --- their nuclei
are, and those nuclei are composed of quarks. Quarks carry fractional charge ($\pm\frac{1}{3}$
and $\pm\frac{2}{3}$ of the elementary charge $e$). They are confined within hadrons by the strong
force (QCD, quantum chromodynamics), mediated by eight massless gluons.

\begin{center}
\rowcolors{2}{rowgray}{white}
\begin{tabularx}{\textwidth}{L{2cm}L{2.2cm}L{2.8cm}XC{1.5cm}}
\toprule
\rowcolor{primary}
\tableheader{Category} & \tableheader{Particle} & \tableheader{Mass (MeV/$c^2$)} & \tableheader{Role} & \tableheader{Charge} \\
\midrule
Quark & Up (u) & 2.2 & Proton/neutron constituent & $+2/3$ \\
Quark & Down (d) & 4.7 & Proton/neutron constituent & $-1/3$ \\
Quark & Strange (s) & 96 & Kaon constituent & $-1/3$ \\
Quark & Charm (c) & 1,270 & D meson constituent & $+2/3$ \\
Quark & Bottom (b) & 4,180 & B meson, LHCb studies & $-1/3$ \\
Quark & Top (t) & 172,690 & Heaviest SM particle & $+2/3$ \\
Lepton & Electron ($e$) & 0.511 & Atomic structure, signals & $-1$ \\
Lepton & Muon ($\mu$) & 105.7 & Penetrating decay product & $-1$ \\
Lepton & Tau ($\tau$) & 1,777 & Heavy lepton & $-1$ \\
Lepton & $\nu_e$, $\nu_\mu$, $\nu_\tau$ & $< 2 \times 10^{-3}$ & Weak interaction & 0 \\
Boson & Photon ($\gamma$) & 0 & EM force carrier & 0 \\
Boson & W$^\pm$ & 80,369 & Weak force carrier (CC) & $\pm 1$ \\
Boson & Z$^0$ & 91,188 & Weak force carrier (NC) & 0 \\
Boson & Gluon ($g$) & 0 & Strong force carrier, 8 types & 0 \\
Boson & Higgs ($H^0$) & 125,100 & Mass generation (BEH mech.) & 0 \\
\bottomrule
\end{tabularx}
\end{center}

\textit{Source: Particle Data Group, Review of Particle Physics (2024), Physical Review D 110, 030001.}

\subsubsection{LHC: Engineering Parameters (Levels 86--90)}

\begin{center}
\rowcolors{2}{rowgray}{white}
\begin{tabularx}{\textwidth}{L{5cm}X}
\toprule
\rowcolor{primary}
\tableheader{Parameter} & \tableheader{Value (Run 3, 2022--2025)} \\
\midrule
Circumference & 26,658.883 m \\
Depth below surface & 45--175 m \\
Beam energy (Run 3) & 6.8 TeV per beam \\
Center-of-mass energy & $\sqrt{s} = 13.6$ TeV \\
Design maximum $\sqrt{s}$ & 14 TeV \\
Main dipole magnets & 1,232 (NbTi superconducting) \\
Dipole magnetic field & 8.60 T (Run 3) \\
Dipole operating temperature & 1.9 K (superfluid He-II) \\
Dipole operating current & 11,850 A \\
Bunch spacing & 25 ns (40.079 MHz crossing rate) \\
Protons per bunch & $\sim 1.15 \times 10^{11}$ \\
Peak luminosity (Run 3) & $2.0 \times 10^{34}$ cm$^{-2}$s$^{-1}$ \\
Average pile-up ($\mu$) & $\sim$45 inelastic pp interactions/crossing \\
Proton velocity & $0.999999991c$ at 6.8 TeV ($\gamma \approx 7253$) \\
Data rate after L1 trigger & $\approx$100 kHz $\rightarrow$ 1--4 kHz to tape \\
\bottomrule
\end{tabularx}
\end{center}

\textit{Source: CERN Document Server, LHC Machine Outreach and Technical Design Reports.}

\subsubsection{Data Acquisition and AI in the High-Level Trigger (Levels 88--92)}

The LHC produces 40 million bunch crossings per second. Each crossing contains $\sim$45
inelastic collisions (pile-up). Raw ATLAS detector data rate: $\approx$60 TB/s (all subdetectors
combined). Hardware Level-1 (L1) trigger reduces to $\sim$100 kHz in 2.5 $\mu$s. Software
High-Level Trigger (HLT) reduces to $\sim$1--4 kHz in $\sim$200 ms per event, using a farm of
$\sim$50,000 CPU/GPU cores.

\textbf{Graph Neural Networks for track reconstruction} (Duarte et al., 2020; Caillou et al.,
2022): Charged particles leave discrete hits in the silicon pixel and strip detectors (ATLAS inner
detector: 86 million channels). Track reconstruction must identify which hits belong to the same
particle and fit a helical track to the hits (in the 2 T solenoid field). Traditional Kalman filter
approaches scale as $O(n^3)$ in hit occupancy; at Run 3 pile-up ($\sim$45), this is a 90,000-hit
pattern recognition problem executed $\sim$100,000 times per second. GNN approach: represent each
hit as a node (features: $r$, $\phi$, $z$, layer ID), each pair of geometrically adjacent hits as a
directed edge, classify edges as belonging to the same track. GNN achieves tracking efficiency
$> 99\%$ at speed compatible with HLT requirements (arXiv:2007.13681). This is AI doing the same
job as the spring terminal: preserving fidelity across a complex interface, 100,000 times per
second.

% ─────────────────────────────────────────────────────────────────
\subsection{M10 --- FIDELITY: Engineering Fidelity as Meta-Layer}
% ─────────────────────────────────────────────────────────────────

\subsubsection{The Calibration Chain: 17 Orders of Magnitude}

Engineering fidelity is the degree to which a transformation preserves the information content
of its input. At every layer of the construction set, fidelity can be defined, measured, and
improved. The remarkable fact is that all ten layers share the same underlying principle ---
and their calibration chains are physically connected through NIST traceability.

\begin{center}
\rowcolors{2}{rowgray}{white}
\begin{tabularx}{\textwidth}{L{1.8cm}L{2.8cm}L{2.5cm}X}
\toprule
\rowcolor{primary}
\tableheader{Layer} & \tableheader{Fidelity Metric} & \tableheader{Units} & \tableheader{Calibration Reference} \\
\midrule
M01 Matter & Contact resistance & m$\Omega$ & NIST SRM resistors \\
M02 Signal & SNR, THD & dB & IEC 60268; NIST 1-$\Omega$ standard \\
M03 Logic & Bit error rate & BER & ETSI TS 136 series \\
M04 Compute & IPC, cache miss \% & IPC, \% & SPEC CPU 2017 benchmarks \\
M05 Software & UB rate, cyclomatic $\sigma$ & UB/KLOC & ISO/IEC 25010 \\
M06 Network & Packet loss, jitter & \%, ms & RFC 3393; IEEE 802.1Qbv \\
M07 Intelligence & Cross-entropy loss & nats & MLPerf Training v4.0 \\
M08 Quantum & Coherence $T_2$, gate fidel. & $\mu$s, \% & NIST quantum benchmarks \\
M09 Collider & Energy res. $\sigma_E/E$ & \% & GEANT4 Monte Carlo + $Z \rightarrow ee$ \\
M10 Meta & Calibration depth, traceability & \# standards in chain & NIST primary standards \\
\bottomrule
\end{tabularx}
\end{center}

\subsubsection{NIST Traceability: The Universal Reference (Levels 95--100)}

Every measurement in engineering ultimately traces to a national metrology institute primary
standard. The spring terminal contact resistance is measured with a milliohm meter calibrated
against a NIST-traceable resistance standard (e.g., NIST SRM 1482, 1484, 1485 for resistance
standards traceable to the quantum Hall effect, which is reproducible to $10^{-9}$ via
$R_K = h/e^2 = 25812.807$ $\Omega$). The ATLAS hadronic calorimeter's energy scale is calibrated
using $Z \rightarrow ee$ events: the Z boson mass $m_Z = 91.1876 \pm 0.0021$ GeV/$c^2$ (PDG 2024)
is known to six significant figures, providing an \textit{in situ} calibration source. The Higgs
boson mass is then measured as $m_H = 125.20 \pm 0.11$ GeV/$c^2$ (ATLAS and CMS combined, 2023).

The calibration chain from the spring terminal (15 cents, $10^{-3}$ $\Omega$ precision) to the
Higgs mass ($125 \times 10^9$ eV, $10^{-3}$ relative precision) spans 17 orders of magnitude in
energy scale and uses the same logical structure at every step: compare your measurement to a known
reference, quantify the uncertainty, propagate it faithfully. This is Level 100. This is what it
means to understand engineering fidelity.

\subsection{Safety and Sensitivity Considerations}

\begin{center}
\rowcolors{2}{rowgray}{white}
\begin{tabularx}{\textwidth}{L{3.5cm}L{1.5cm}X}
\toprule
\rowcolor{primary}
\tableheader{Consideration} & \tableheader{Layer} & \tableheader{Mitigation} \\
\midrule
Electrical safety (mains) & M02 & All L1--20 experiments at $\leq$9 V; 120/230 V mains explicitly out of scope \\
Radioactive element samples & M01 & Data-only coverage; no radioactive or CBRN-relevant element samples \\
Dual-use RF power & M06 & Curriculum specifies low-power only; regulatory refs. noted \\
AI bias and dataset fidelity & M07 & Bias framed as distributional mismatch; fidelity language used \\
Particle accelerator access & M09 & All LHC content is observational/theoretical; no operational access implied \\
Nuclear/weapons dual-use & M09 & Standard Model = basic science; no weapons-applicable parameters \\
Element synthesis (superheavy) & M01 & Oganesson through Nihonium covered data-only; no synthesis details \\
\bottomrule
\end{tabularx}
\end{center}

\subsection{Source Bibliography}

\subsubsection{Government and Agency Sources}

\begin{itemize}[leftmargin=*]
  \item CERN. (2022--2025). \textit{LHC Run 3 Technical Documentation and Machine Parameters}. CERN Document Server. \texttt{cds.cern.ch}
  \item NIST. (2024). \textit{Atomic Weights and Isotopic Compositions with Relative Atomic Masses}. \texttt{physics.nist.gov/comp}
  \item NIST. (2018). \textit{CODATA Recommended Values of Fundamental Physical Constants}. \texttt{physics.nist.gov/constants}
  \item NIST. (2024). \textit{Standard Reference Materials Catalog --- Electrical Resistance Standards}. \texttt{nist.gov/srm}
  \item Particle Data Group. (2024). \textit{Review of Particle Physics}. Physical Review D, 110, 030001. \texttt{pdg.lbl.gov}
  \item IEEE Std 802.11-2020. \textit{IEEE Standard for Information Technology --- Wireless LAN MAC and PHY Specifications}. IEEE Standards Association.
  \item IEEE Std 802.1Qbv-2015. \textit{Enhancements for Scheduled Traffic (Time-Sensitive Networking)}. IEEE Standards Association.
  \item IEC 60268-4. \textit{Sound System Equipment --- Part 4: Microphones}. International Electrotechnical Commission.
\end{itemize}

\subsubsection{Peer-Reviewed Research}

\begin{itemize}[leftmargin=*]
  \item Rosenblatt, F. (1958). The Perceptron: A Probabilistic Model for Information Storage and Organization in the Brain. \textit{Psychological Review}, 65(6), 386--408.
  \item Shannon, C. E. (1948). A Mathematical Theory of Communication. \textit{Bell System Technical Journal}, 27(3), 379--423.
  \item He, K., Zhang, X., Ren, S., \& Sun, J. (2016). Deep Residual Learning for Image Recognition. \textit{IEEE CVPR}. arXiv:1512.03385.
  \item Vaswani, A. et al. (2017). Attention Is All You Need. \textit{NeurIPS 2017}. arXiv:1706.03762.
  \item Duarte, J. et al. (2020). Graph Neural Networks for Particle Physics. \textit{Machine Learning: Science and Technology}, 1(4). arXiv:2007.13681.
  \item Arndt, M. et al. (1999). Wave-particle duality of $C_{60}$ molecules. \textit{Nature}, 401, 680--682.
  \item Aspect, A., Grangier, P., \& Roger, G. (1982). Experimental Realization of Einstein-Podolsky-Rosen-Bohm Gedankenexperiment. \textit{Physical Review Letters}, 49(2), 91--94.
  \item von Neumann, J. (1945). \textit{First Draft of a Report on the EDVAC}. University of Pennsylvania Technical Report. Moore School of Electrical Engineering.
  \item Hennessy, J. L., \& Patterson, D. A. (1990). \textit{Computer Architecture: A Quantitative Approach} (1st ed.). Morgan Kaufmann.
\end{itemize}

\subsubsection{Professional Organizations and Curriculum Sources}

\begin{itemize}[leftmargin=*]
  \item ACM. (2023). \textit{Computer Science Curricula 2023}. Association for Computing Machinery. \texttt{cs2023.acm.org}
  \item IEEE Computer Society. (2024). \textit{Guide to the Software Engineering Body of Knowledge (SWEBOK 4.0)}.
  \item Royal Society of Chemistry. \textit{Periodic Table with data}. \texttt{rsc.org/periodic-table}
  \item NVIDIA. (2025). \textit{RTX 5090 Product Brief and GPU Architecture White Paper}. \texttt{nvidia.com/rtx}
  \item MIT OpenCourseWare. (2024). 6.004 Computation Structures; 6.012 Microelectronic Devices and Circuits; 8.05 Quantum Physics II. \texttt{ocw.mit.edu}
  \item IBM Quantum. (2024). \textit{IBM Quantum System Two and Heron Processor Technical Specifications}. \texttt{quantum.ibm.com}
  \item Google Quantum AI. (2024). \textit{Willow: Achieving Below-Threshold Quantum Error Correction}. \texttt{quantumai.google}
\end{itemize}

\newpage

% ==================================================================
%  STAGE 3 — MISSION PACKAGE
% ==================================================================
\stageheader{STAGE 3 --- MISSION PACKAGE}{tertiary}

\section{Milestone Specification}

\metafield{Milestone ID:}{CONSTRUCTION-SET-v1.0}
\metafield{Activation Profile:}{Fleet (Full Crew, 20 roles)}
\metafield{Parallel Tracks:}{3 (Wave 1A, 1B, 1C)}
\metafield{Wave Count:}{6 (W0, W1A, W1B, W1C, W2, W3)}
\metafield{Model Budget:}{Opus 30\% / Sonnet 60\% / Haiku 10\%}
\metafield{Token Budget:}{160,000 total}

\subsection{Mission Objective}

Produce a publication-ready 100-level engineering curriculum reference spanning ten modules from
the periodic table of elements (M01-MATTER, Levels 1--10) through LHC particle physics
(M09-COLLIDER, Levels 81--90) to engineering fidelity as universal meta-principle (M10-FIDELITY,
Levels 91--100), structured for GSD skill-creator delivery. The curriculum must be self-contained
for a high-school physics background, technically rigorous enough to support a physics PhD student,
and architecturally decomposed so that each of the 100 levels is a potential skill-creator mission
component with defined input, output, and fidelity metric.

\subsection{Architecture Overview}

\begin{asciibox}
CONSTRUCTION SET v1.0 --- WAVE EXECUTION ARCHITECTURE

Wave 0  [Sequential, Haiku, 6K tokens]
  Schema definition: fidelity metrics, level taxonomy, source index, LaTeX macros
  GATE: Schema validates against all module draft structures before Wave 1

Wave 1A [Parallel Track A, Sonnet, 32K tokens]
  M01-MATTER + M02-SIGNAL + M03-LOGIC (Levels 1--30)

Wave 1B [Parallel Track B, Sonnet, 30K tokens]      } PARALLEL
  M04-COMPUTE + M05-SOFTWARE + M06-NETWORK (L31--60)  } across 1A,
                                                       } 1B, 1C
Wave 1C [Parallel Track C, Opus, 60K tokens]        }
  M07-INTELLIGENCE + M08-QUANTUM + M09-COLLIDER + M10-FIDELITY (L61--100)

Wave 2  [Sequential, Opus, 20K tokens]
  Layer bridges (9x) + unified calibration chain + Standard Model-periodic table link
  GATE: Human review (CAPCOM) before Wave 3

Wave 3  [Sequential, Sonnet, 12K tokens]
  Source citation audit + level numbering + XeLaTeX compile (3 passes) + index
  GATE: CAPCOM final sign-off
  OUTPUT: D13 Full Curriculum PDF
\end{asciibox}

\subsection{System Layers}

\begin{enumerate}[leftmargin=*, label=\textbf{SL-\arabic*.}]
  \item \textbf{Foundation Layer} (W0): Fidelity metrics schema (JSON), level taxonomy (JSON: levels 1--100 with module and prerequisite), source citation index, shared LaTeX macros.
  \item \textbf{Foundational Physics Layer} (W1A): M01 Matter, M02 Signal, M03 Logic --- levels 1--30.
  \item \textbf{Computing Systems Layer} (W1B): M04 Compute, M05 Software, M06 Network --- levels 31--60.
  \item \textbf{Advanced Domains Layer} (W1C): M07 Intelligence, M08 Quantum, M09 Collider, M10 Fidelity --- levels 61--100.
  \item \textbf{Integration Layer} (W2): Cross-module bridges, unified calibration chain, narrative synthesis.
  \item \textbf{Publication Layer} (W3): Verification, LaTeX compilation, indexing, final PDF.
\end{enumerate}

\subsection{Deliverables}

\begin{center}
\rowcolors{2}{rowgray}{white}
\begin{tabularx}{\textwidth}{C{0.8cm}L{3.5cm}L{3.5cm}X}
\toprule
\rowcolor{primary}
\tableheader{\#} & \tableheader{Deliverable} & \tableheader{Acceptance Criterion} & \tableheader{Format} \\
\midrule
D1 & Shared Schema & JSON validates; all 10 modules reference it & JSON + LaTeX macros \\
D2 & M01 MATTER & 40+ elements mapped, spring terminal BOM & PDF section + BOM table \\
D3 & M02 SIGNAL & L11--20 experiment sequence complete & PDF section + BOM \\
D4 & M03 LOGIC & 7 gates + CMOS + IEEE754 + bfloat16 & PDF section \\
D5 & M04 COMPUTE & MCU to superscalar + Amiga case study & PDF section \\
D6 & M05 SOFTWARE & Compiler pipeline + OS + Rust safety & PDF section \\
D7 & M06 NETWORK & OSI table + mesh + LHC timing network & PDF section \\
D8 & M07 INTELLIGENCE & Perceptron $\rightarrow$ transformer + GPU table & PDF section \\
D9 & M08 QUANTUM & QM foundations + coherence table & PDF section \\
D10 & M09 COLLIDER & LHC parameters + SM table + GNN HLT & PDF section \\
D11 & M10 FIDELITY & Calibration chain table, all 10 layers & PDF section \\
D12 & Integration & 9 layer bridges + unified chain + narrative & PDF section \\
D13 & Full Curriculum PDF & All 100 levels indexed, compiles clean & Single XeLaTeX PDF \\
\bottomrule
\end{tabularx}
\end{center}

\subsection{Component Breakdown}

\begin{center}
\rowcolors{2}{rowgray}{white}
\begin{tabularx}{\textwidth}{L{3cm}L{3cm}C{1.6cm}C{1.6cm}}
\toprule
\rowcolor{primary}
\tableheader{Component} & \tableheader{Dependencies} & \tableheader{Model} & \tableheader{Tokens} \\
\midrule
W0: Schema \& Index & None & Haiku & 6,000 \\
W1A: M01 MATTER & Schema & Sonnet & 12,000 \\
W1A: M02 SIGNAL & M01 & Sonnet & 10,000 \\
W1A: M03 LOGIC & M02 & Sonnet & 10,000 \\
W1B: M04 COMPUTE & M03, Schema & Sonnet & 12,000 \\
W1B: M05 SOFTWARE & M04 & Sonnet & 10,000 \\
W1B: M06 NETWORK & M04, M05 & Sonnet & 8,000 \\
W1C: M07 INTELLIGENCE & M04, M06, Schema & Opus & 16,000 \\
W1C: M08 QUANTUM & M01, M03, Schema & Opus & 14,000 \\
W1C: M09 COLLIDER & M07, M08, Schema & Opus & 16,000 \\
W1C: M10 FIDELITY & All M01--M09 & Opus & 14,000 \\
W2: Integration & All Wave 1 & Opus & 20,000 \\
W3: Publication & Wave 2 & Sonnet & 12,000 \\
\midrule
\multicolumn{3}{r}{\textbf{Total}} & \textbf{160,000} \\
\bottomrule
\end{tabularx}
\end{center}

\subsection{Model Assignment Rationale}

\begin{center}
\rowcolors{2}{rowgray}{white}
\begin{tabularx}{\textwidth}{L{1.8cm}C{1.4cm}XC{1.5cm}}
\toprule
\rowcolor{primary}
\tableheader{Model} & \tableheader{Share} & \tableheader{Rationale} & \tableheader{Tokens} \\
\midrule
Opus & 30\% & Cross-domain synthesis, Standard Model, QM, AI-hardware bridge, judgment-heavy integration & 48,000 \\
Sonnet & 60\% & Document production, structured content, technical surveys, publication assembly & 96,000 \\
Haiku & 10\% & Schema scaffolding, level taxonomy, template generation & 16,000 \\
\bottomrule
\end{tabularx}
\end{center}

\newpage

\section{Wave Execution Plan}

\subsection{Wave Summary}

\begin{center}
\rowcolors{2}{rowgray}{white}
\begin{tabularx}{\textwidth}{L{1.6cm}L{1.8cm}L{3.5cm}L{1.8cm}X}
\toprule
\rowcolor{primary}
\tableheader{Wave} & \tableheader{Mode} & \tableheader{Components} & \tableheader{Model} & \tableheader{Gate Condition} \\
\midrule
W0 & Sequential & Schema, taxonomy, index & Haiku & JSON schema validates \\
W1A & Parallel & M01, M02, M03 & Sonnet & 3 module drafts complete \\
W1B & Parallel & M04, M05, M06 & Sonnet & 3 module drafts complete \\
W1C & Parallel & M07, M08, M09, M10 & Opus & 4 module drafts complete \\
W2 & Sequential & Integration, bridges & Opus & Human review (CAPCOM) \\
W3 & Sequential & Compile, verify, index & Sonnet & CAPCOM final sign-off \\
\bottomrule
\end{tabularx}
\end{center}

\subsection{Wave 0: Foundation}

\begin{center}
\rowcolors{2}{rowgray}{white}
\begin{tabularx}{\textwidth}{L{2.8cm}XL{2cm}}
\toprule
\rowcolor{primary}
\tableheader{Task} & \tableheader{Output} & \tableheader{Agent} \\
\midrule
W0.1 Fidelity schema & JSON: \{layer, metric, unit, reference, calibration\_std\} per 10 layers & Haiku \\
W0.2 Level taxonomy & JSON: level 1..100, module, prerequisites, fidelity\_metric & Haiku \\
W0.3 Source index & Deduplicated bibliography seeds; DOI or URL per source & Haiku \\
W0.4 LaTeX macros & Shared environments; fidelity table macros; layer color codes & Haiku \\
\bottomrule
\end{tabularx}
\end{center}

\subsection{Wave 1: Parallel Research Tracks}

\textbf{Wave 1A (Sonnet) --- Matter, Signal, Logic:}
\begin{center}
\rowcolors{2}{rowgray}{white}
\begin{tabularx}{\textwidth}{L{2.2cm}XL{2.2cm}}
\toprule
\rowcolor{primary}
\tableheader{Task} & \tableheader{Output} & \tableheader{Agent} \\
\midrule
W1A.1 Element table & 40+ elements: Z, symbol, $\rho$/$E_g$/role, device example & CRAFT-MATTER \\
W1A.2 Spring terminal & Physical analysis, element composition, BOM, contact resistance spec & CRAFT-MATTER \\
W1A.3 L11--20 sequence & Step-by-step experiment per level, BOM, measurable outcome & CRAFT-SIGNAL \\
W1A.4 RC/555 section & Reactive components, resonance, 555 timer in astable mode & CRAFT-SIGNAL \\
W1A.5 Logic gates & 7 gates, truth tables, CMOS implementation, functional completeness & CRAFT-LOGIC \\
W1A.6 Arithmetic + FP & Binary, 2's complement, IEEE 754, bfloat16, machine epsilon table & CRAFT-LOGIC \\
\bottomrule
\end{tabularx}
\end{center}

\textbf{Wave 1B (Sonnet) --- Compute, Software, Network:}
\begin{center}
\rowcolors{2}{rowgray}{white}
\begin{tabularx}{\textwidth}{L{2.2cm}XL{2.2cm}}
\toprule
\rowcolor{primary}
\tableheader{Task} & \tableheader{Output} & \tableheader{Agent} \\
\midrule
W1B.1 MCU curriculum & ATmega328P spec, L31--L40 experiment sequence, ADC calibration loop & CRAFT-COMPUTE \\
W1B.2 CPU architecture & 5-stage pipeline, OoO superscalar, Amiga 68000 chipset case study & CRAFT-COMPUTE \\
W1B.3 Memory table & L1/L2/L3/RAM/SSD/HDD latency and capacity table (2025 values) & CRAFT-COMPUTE \\
W1B.4 Compiler pipeline & LLVM passes, UB taxonomy, Rust ownership model & CRAFT-SOFT \\
W1B.5 OS section & Kernel, syscall, virtual memory, TLB, POSIX portability & CRAFT-SOFT \\
W1B.6 OSI table & 7-layer fidelity analysis with protocols & CRAFT-NET \\
W1B.7 Mesh + LHC timing & GSD mesh spec + LHC TTC synchronization comparison & CRAFT-NET \\
\bottomrule
\end{tabularx}
\end{center}

\textbf{Wave 1C (Opus) --- Intelligence, Quantum, Collider, Fidelity:}
\begin{center}
\rowcolors{2}{rowgray}{white}
\begin{tabularx}{\textwidth}{L{2.2cm}XL{2.2cm}}
\toprule
\rowcolor{primary}
\tableheader{Task} & \tableheader{Output} & \tableheader{Agent} \\
\midrule
W1C.1 ML ladder & Perceptron to transformer: hardware substrate table per milestone & CRAFT-AI \\
W1C.2 GPU section & RTX 5090 specs, CUDA, bfloat16 rationale, TFLOPS vs CPU & CRAFT-AI \\
W1C.3 Fidelity-loss link & Cross-entropy, gradient vanishing, ResNet residuals as fidelity eng. & CRAFT-AI \\
W1C.4 Double slit & Fidelity failure at quantum measurement; wave function, Born rule & CRAFT-QUANT \\
W1C.5 Qubit coherence & $T_2$ table (6 platforms), error correction as redundancy fidelity & CRAFT-QUANT \\
W1C.6 SM particle table & 15 particles: mass (MeV/$c^2$), charge, role, source PDG 2024 & CRAFT-PHYS \\
W1C.7 LHC parameter table & All Run 3 parameters with CERN source attribution & CRAFT-PHYS \\
W1C.8 GNN HLT section & Track reconstruction problem, GNN approach, efficiency metric & CRAFT-PHYS \\
W1C.9 Calibration chain & 10-layer fidelity table + NIST traceability + spring-to-Higgs narrative & CRAFT-FIDEL \\
\bottomrule
\end{tabularx}
\end{center}

\subsection{Wave 2: Synthesis}

\begin{center}
\rowcolors{2}{rowgray}{white}
\begin{tabularx}{\textwidth}{L{2.5cm}XL{2cm}}
\toprule
\rowcolor{primary}
\tableheader{Task} & \tableheader{Output} & \tableheader{Agent} \\
\midrule
W2.1 Layer bridges & 9 explicit transition documents (L10$\rightarrow$11 through L90$\rightarrow$91) with bridging concept & INTEG \\
W2.2 Calibration chain & Single master table: all 10 layers, metric, unit, NIST link & INTEG \\
W2.3 SM $\leftrightarrow$ PT link & Narrative: periodic table $\rightarrow$ Standard Model as successor construction set & Opus \\
W2.4 HW-AI-LHC bridge & Spring terminal $\rightarrow$ GPU Tensor Core $\rightarrow$ LHC GNN trigger chain & Opus \\
W2.5 Level index & 100-entry index: level number, module, title, fidelity metric & EXEC-2 \\
W2.6 Human review & CAPCOM presents integrated document; Tibsfox review gate & CAPCOM \\
\bottomrule
\end{tabularx}
\end{center}

\subsection{Wave 3: Publication and Verification}

\begin{center}
\rowcolors{2}{rowgray}{white}
\begin{tabularx}{\textwidth}{L{2.5cm}XL{2cm}}
\toprule
\rowcolor{primary}
\tableheader{Task} & \tableheader{Output} & \tableheader{Agent} \\
\midrule
W3.1 Citation audit & SC-SRC test: all citations verified against primary sources & VERIFY \\
W3.2 Numerical audit & SC-NUM test: all quantitative claims attributed & VERIFY \\
W3.3 Level numbering & All 100 levels present; no gaps; all prerequisites valid & VERIFY \\
W3.4 Safety checks & SC-EXP: all L1--20 experiments $\leq$9V; SC-DU: no dual-use content & VERIFY \\
W3.5 XeLaTeX compile & Three-pass compile; zero errors and zero unresolved references & EXEC-1 \\
W3.6 Final PDF & D13: publication-ready curriculum, TOC, level index & EXEC-1 \\
W3.7 Sign-off & CAPCOM final gate; LOG records commit message & CAPCOM \\
\bottomrule
\end{tabularx}
\end{center}

\subsection{Cache Optimization Strategy}

\begin{itemize}[leftmargin=*]
  \item The shared schema (W0: 6K tokens) is sized as a 5-minute cache-warm target. All Wave 1 agents receive it at session start and exploit the 300-second TTL window. After the first module call, the schema header is cached, costing zero marginal tokens for the remaining 9 module calls within the same session.
  \item Wave 1A, 1B, and 1C are each sized to complete within a single 200K-token context window per track (32K, 30K, 60K respectively), enabling within-session caching of prior module outputs for subsequent modules on the same track.
  \item Wave 2 receives all 10 module outputs as a single structured document (not 10 separate files), minimizing context assembly overhead. The integration agent exploits schema consistency to navigate the full reference without reloading.
  \item Wave 3 receives the Wave 2 integration document as its primary input; the citation audit operates on the embedded bibliography, not the full source reference, keeping context window overhead minimal.
\end{itemize}

\subsection{Token Budget}

\begin{center}
\rowcolors{2}{rowgray}{white}
\begin{tabularx}{\textwidth}{L{3cm}C{2.5cm}C{2cm}C{2cm}}
\toprule
\rowcolor{primary}
\tableheader{Wave} & \tableheader{Context Windows} & \tableheader{Model} & \tableheader{Tokens} \\
\midrule
Wave 0 & 1 sequential & Haiku & 6,000 \\
Wave 1A & 3 parallel & Sonnet & 32,000 \\
Wave 1B & 3 parallel & Sonnet & 30,000 \\
Wave 1C & 4 parallel & Opus & 60,000 \\
Wave 2 & 2 sequential & Opus & 20,000 \\
Wave 3 & 2 sequential & Sonnet & 12,000 \\
\midrule
\textbf{Total} & \textbf{$\sim$15 windows} & & \textbf{160,000} \\
\bottomrule
\end{tabularx}
\end{center}

\subsection{Dependency Graph}

\begin{asciibox}
W0: Schema + Level Taxonomy + Source Index
 |
 +---[1A]---> M01-MATTER
 |                 |
 |            M02-SIGNAL  <-- M01
 |                 |
 |            M03-LOGIC   <-- M02
 |                 |
 +---[1B]---> M04-COMPUTE <-- M03, Schema
 |                 |
 |            M05-SOFTWARE <-- M04
 |                 |
 |            M06-NETWORK  <-- M04, M05
 |                 |
 +---[1C]---> M07-INTELLIGENCE <-- M04, M06, Schema
 |            M08-QUANTUM      <-- M01, M03, Schema
 |            M09-COLLIDER     <-- M07, M08, Schema
 |            M10-FIDELITY     <-- ALL M01-M09
 |
W2: Integration <-- All Wave 1 outputs
 |   (9 bridges + calibration chain + SM-PT narrative)
 |   HUMAN GATE: Tibsfox review
 |
W3: Publication <-- W2 Integration
     (Verify + Compile + Index)
     CAPCOM GATE: Final sign-off
     |
     D13: The Infinite and One Construction Set
          Full 100-Level Curriculum PDF
\end{asciibox}

\newpage

\section{Test Plan}

\subsection{Test Categories}

\begin{center}
\rowcolors{2}{rowgray}{white}
\begin{tabularx}{\textwidth}{L{3cm}XC{1.8cm}C{1.6cm}C{1.3cm}}
\toprule
\rowcolor{primary}
\tableheader{Category} & \tableheader{Purpose} & \tableheader{Priority} & \tableheader{On Failure} & \tableheader{Count} \\
\midrule
Safety-critical & Non-negotiable accuracy and safety boundaries & Mandatory & BLOCK & 6 \\
Core functionality & Deliverable acceptance criteria & Required & BLOCK & 24 \\
Integration & Cross-module and cross-layer validation & Required & BLOCK & 10 \\
Edge cases & Robustness and completeness & Best-effort & LOG & 7 \\
\midrule
\multicolumn{4}{r}{\textbf{Total}} & \textbf{47} \\
\bottomrule
\end{tabularx}
\end{center}

\subsection{Safety-Critical Tests}

\begin{center}
\rowcolors{2}{rowgray}{white}
\begin{tabularx}{\textwidth}{C{1.5cm}L{3.5cm}X}
\toprule
\rowcolor{primary}
\tableheader{Test ID} & \tableheader{Verifies} & \tableheader{Expected Behavior --- BLOCK on failure} \\
\midrule
SC-SRC & Source quality & All citations traceable to CERN, NIST, IEEE, PDG, ACM, peer-reviewed journals. Zero entertainment media, blogs, or unsourced claims. BLOCK. \\
SC-NUM & Numerical attribution & Every physical constant, detector spec, transistor count, GPU spec, energy value attributed to specific primary source. BLOCK. \\
SC-ADV & No policy advocacy & Document presents physics and engineering evidence without advocating for funding positions, educational policy, or institutional agendas. BLOCK. \\
SC-EXP & Experiment safety & All L1--L20 hands-on experiments specified at $\leq$9 V supply. No mains-voltage experiments anywhere in scope. BLOCK. \\
SC-DU & Dual-use check & No operational accelerator control details; Standard Model content is unambiguous basic science education; no weapons-applicable nuclear or radiological parameters. BLOCK. \\
SC-ELE & Element safety & No radioactive, highly toxic, or CBRN-relevant element samples specified for student hands-on handling. Data-only coverage for such elements. BLOCK. \\
\bottomrule
\end{tabularx}
\end{center}

\subsection{Core Functionality Tests (Selected)}

\begin{center}
\rowcolors{2}{rowgray}{white}
\begin{tabularx}{\textwidth}{C{1.5cm}L{3.8cm}X}
\toprule
\rowcolor{primary}
\tableheader{Test ID} & \tableheader{Verifies} & \tableheader{Expected Outcome} \\
\midrule
CF-01 & Periodic table completeness & $\geq$40 elements with Z, symbol, primary engineering application, and device example. All sourced. \\
CF-02 & L1--20 experiment sequence & Complete BOM and step-by-step procedure for each of 20 levels; measurable outcome stated per level. \\
CF-03 & Layer transition bridges & 9 explicit bridges (L10$\rightarrow$11 through L90$\rightarrow$91) each containing a named bridging concept. \\
CF-04 & Amiga case study present & Amiga 68000 chipset architecture documented with Agnus/Denise/Paula roles and Amiga Principle lesson. \\
CF-05 & GPU specifications & RTX 5090 (or current-gen 2025): TFLOPS FP32, bfloat16 Tensor TFLOPS, VRAM, bandwidth cited to NVIDIA. \\
CF-06 & LHC parameter table & All 15 listed Run 3 parameters present with CERN attribution. \\
CF-07 & Standard Model table & $\geq$15 particles with mass (MeV/$c^2$), charge, role. All sourced to PDG 2024. \\
CF-08 & Qubit coherence table & $T_2$ values for $\geq$4 hardware platforms with year and technology column. \\
CF-09 & Calibration chain table & All 10 layers present with fidelity metric, unit, and named calibration standard. \\
CF-10 & Level index & All 100 levels 1--100 present in index; no gaps; each assigned to a module. \\
CF-11 & bfloat16 rationale & bfloat16 versus float32 comparison: exponent bits, mantissa bits, AI training rationale documented. \\
CF-12 & GNN tracking section & Graph Neural Network for LHC track reconstruction: problem statement, approach, efficiency metric. \\
\bottomrule
\end{tabularx}
\end{center}

\subsection{Integration Tests (Selected)}

\begin{center}
\rowcolors{2}{rowgray}{white}
\begin{tabularx}{\textwidth}{C{1.5cm}L{3.8cm}X}
\toprule
\rowcolor{primary}
\tableheader{Test ID} & \tableheader{Verifies} & \tableheader{Expected Outcome} \\
\midrule
IN-01 & M01 $\rightarrow$ M09 (Nb/Ti link) & Niobium and titanium in M01 element table are explicitly cross-referenced to LHC dipole magnets in M09. \\
IN-02 & M03 $\rightarrow$ M07 (bfloat16) & bfloat16 format introduced in M03 (floating point) is cross-referenced to AI training implications in M07. \\
IN-03 & M04 $\rightarrow$ M07 (GPU vs CPU) & CPU architecture from M04 provides the baseline comparison for GPU parallelism in M07. \\
IN-04 & M06 $\rightarrow$ M09 (timing) & LHC TTC timing network in M09 is explicitly connected to IEEE 802.1Qbv TSN introduced in M06. \\
IN-05 & M10 (master chain) & M10 calibration chain references all 9 prior modules; no module left without a fidelity metric entry. \\
IN-06 & M02 $\rightarrow$ M09 (impedance) & 50 $\Omega$ impedance matching introduced in M02 (L20) explicitly connected to LHC signal cable termination in M09. \\
IN-07 & Self-containment check & A reader with HS physics background can follow from L1 to L100 without encountering undefined terms. \\
IN-08 & Amiga Principle thread & Amiga Principle mentioned in M04 (CRAFT-COMPUTE) is explicitly echoed in at least M07 (AI) and M09 (LHC). \\
\bottomrule
\end{tabularx}
\end{center}

\subsection{Verification Matrix}

\begin{center}
\rowcolors{2}{rowgray}{white}
\begin{tabularx}{\textwidth}{XL{3.5cm}C{1.5cm}}
\toprule
\rowcolor{primary}
\tableheader{Success Criterion} & \tableheader{Test ID(s)} & \tableheader{Status} \\
\midrule
SC-1: 40+ elements with applications & CF-01, SC-SRC, SC-NUM & Pending \\
SC-2: L1--20 experiment sequence & CF-02, SC-EXP, SC-ELE & Pending \\
SC-3: 9 layer transition bridges & CF-03, IN-01 & Pending \\
SC-4: CPU to Amiga to Amiga Principle & CF-04, IN-08 & Pending \\
SC-5: Perceptron to transformer + hardware & CF-05, IN-03 & Pending \\
SC-6: LHC DAQ chain + AI in HLT & CF-06, CF-12, IN-04 & Pending \\
SC-7: 10-layer fidelity calibration chain & CF-09, IN-05 & Pending \\
SC-8: QM bridge classical to quantum & CF-08, IN-02 & Pending \\
SC-9: All quantitative claims attributed & SC-NUM, SC-SRC & Pending \\
SC-10: Self-contained HS physics baseline & IN-07 & Pending \\
SC-11: Three non-adjacent layer links & IN-01, IN-02, IN-06 & Pending \\
SC-12: GSD skill-creator decomposition ready & CF-10, CF-03 & Pending \\
\bottomrule
\end{tabularx}
\end{center}

\section{Mission Crew Manifest}

\begin{center}
\rowcolors{2}{rowgray}{white}
\begin{tabularx}{\textwidth}{L{2.2cm}L{3.2cm}X}
\toprule
\rowcolor{primary}
\tableheader{Role} & \tableheader{Agent} & \tableheader{Responsibility} \\
\midrule
FLIGHT & Orchestrator & Mission direction; go/no-go at wave gates; model allocation decisions \\
PLAN & Planner & Wave decomposition; parallel track optimization; cache strategy \\
SCOUT & Research Agent & Source gathering; CERN/NIST/PDG/IEEE evidence compilation; gap-fill searches \\
EXEC-1 & Document Builder A & Wave 1A production; Wave 3 LaTeX compilation and index \\
EXEC-2 & Document Builder B & Wave 1B production; Wave 2 level index assembly \\
CRAFT-MATTER & Materials Specialist & M01: Periodic table, element applications, spring terminal anatomy \\
CRAFT-SIGNAL & Electrical Specialist & M02: KVL/KCL, experiment sequence L11--20, RC circuits, 555 timer \\
CRAFT-LOGIC & Digital Specialist & M03: Boolean algebra, CMOS, binary arithmetic, IEEE 754, bfloat16 \\
CRAFT-COMPUTE & Architecture Specialist & M04: ATmega328P, CPU pipeline, Amiga 68000, memory hierarchy \\
CRAFT-SOFT & Software Specialist & M05: LLVM pipeline, OS abstraction, Rust ownership, UB taxonomy \\
CRAFT-NET & Network Specialist & M06: OSI model, TCP/IP, mesh topology, LHC TTC timing network \\
CRAFT-AI & Intelligence Specialist & M07: Perceptron to transformer, GPU architecture, gradient fidelity \\
CRAFT-QUANT & Quantum Specialist & M08: Double slit, Schr\"{o}dinger, qubit coherence, error correction \\
CRAFT-PHYS & Collider Specialist & M09: Standard Model, LHC Run 3 parameters, GNN track reconstruction \\
CRAFT-FIDEL & Fidelity Specialist & M10: Calibration chain, NIST traceability, spring-to-Higgs narrative \\
VERIFY & Verification Agent & SC-SRC, SC-NUM, SC-EXP, SC-DU; all 47 tests executed \\
INTEG & Integration Agent & Wave 2 bridge documents; cross-module consistency; unified chain \\
CAPCOM & HITL Interface & Human-machine boundary; Wave 2 review; Wave 3 final sign-off \\
BUDGET & Resource Manager & 160K token budget tracking; Opus/Sonnet/Haiku split enforcement \\
SURGEON & Health Monitor & Research quality degradation detection; source drift alerts \\
LOG & Chronicle Agent & Commit messages; wave summaries; milestone audit trail \\
\bottomrule
\end{tabularx}
\end{center}

\section{Execution Summary}

\begin{center}
\rowcolors{2}{rowgray}{white}
\begin{tabularx}{\textwidth}{L{5.5cm}X}
\toprule
\rowcolor{primary}
\tableheader{Metric} & \tableheader{Value} \\
\midrule
Activation Profile & Fleet (Full Crew, 20 roles) \\
Research Modules & 10 (M01--M10) \\
Curriculum Levels & 100 (L1--L100) \\
Parallel Tracks (Wave 1) & 3 (Wave 1A, 1B, 1C) \\
Wave Count & 6 (W0, W1A, W1B, W1C, W2, W3) \\
Deliverables & 13 (D1--D13) \\
Total Token Budget & 160,000 tokens \\
Estimated Context Windows & $\sim$15 across 6--8 sessions \\
Model Split & Opus 30\% / Sonnet 60\% / Haiku 10\% \\
Human Review Gates & 2 (Wave 2 integration review; Wave 3 final sign-off) \\
Safety-Critical Tests & 6 (all BLOCK on failure) \\
Total Tests & 47 \\
Primary Domains Covered & Matter, Signal, Logic, Compute, Software, Network, AI, Quantum, Particle Physics, Fidelity \\
Physical Scale Range & 15 cents (spring terminal) $\rightarrow$ \$4.7B (LHC construction cost) \\
Energy Scale Range & $10^{-3}$ $\Omega$ contact resistance $\rightarrow$ 125 GeV Higgs mass \\
Key Source Organizations & CERN, NIST, PDG/LBNL, IEEE, ACM, RSC, MIT OCW, NVIDIA, IBM Quantum \\
\bottomrule
\end{tabularx}
\end{center}

\vspace{2em}

\begin{quotebox}
\textit{``A spring terminal on a piece of cardboard. A superconducting dipole at 1.9 Kelvin in a
27-kilometer ring. Between them: the same idea, expressed at different scales --- a connection
made precise, held by a force calibrated to preserve information across a material interface.
The spring constant of the one and the Tesla rating of the other are measurements on the same
calibration chain, separated by seventeen orders of magnitude in energy but united by a single
principle: engineering fidelity. This curriculum makes that chain visible, walkable, and
inexhaustible.''}

\vspace{0.8em}
\textit{--- The Infinite and One Construction Set, through-line}

\vspace{1em}
\textbf{For the GSD ecosystem:} Every layer of this construction set is a demonstration of the
Amiga Principle --- that what a 7 MHz Amiga 500 achieved in 512 KB was not a miracle of raw
power but a miracle of knowing precisely what each chip was for. Agnus, Denise, Paula, and the
68000 each did one thing with perfect fidelity; the system's emergent capability was the product
of their architecture, not the sum of their clock cycles. From the electron to the quark, this
is the lesson. One thread. Ten octaves. Infinite depth.

\vspace{0.8em}
\textit{From spring terminal to Large Hadron Collider: the spaces between are where the meaning lives.}
\end{quotebox}

\end{document}
